% ============================================================================
%  SECTION 1 — INTRODUCTION & MOTIVATION        (Membre 1)
% ============================================================================
\section{Introduction et motivation}
\label{sec:intro}

\subsection{Du cours a une question laissee ouverte}

Le cours d'optimisation convexe etablit un resultat fondamental : pour une
fonction convexe sur un ensemble convexe, \emph{tout minimum local est un
minimum global}. Cette propriete --- que nous appellerons la \guillemotleft{}~propriete-or~\guillemotright{}
--- est ce qui rend les problemes convexes \guillemotleft{}~faciles~\guillemotright{} a optimiser : il
n'existe pas de creux sous-optimal dans lequel un algorithme pourrait rester
piege.

Le cours a ensuite presente deux familles de methodes de resolution :
\begin{itemize}
  \item pour un objectif \textbf{lisse} (differentiable), les methodes fondees
        sur le gradient (descente de gradient, methode de Newton) ;
  \item pour un objectif \textbf{lineaire}, l'algorithme du simplexe.
\end{itemize}

Ces deux familles couvrent une large part des problemes, mais elles partagent
une hypothese implicite forte : la descente de gradient et la methode de Newton
exigent que la fonction objectif soit \emph{differentiable}, tandis que le
simplexe ne traite que le cas \emph{lineaire}. Une question reste donc ouverte,
et c'est precisement celle qu'explore ce projet :

\begin{center}
\itshape
Comment resoudre un probleme convexe lorsque la fonction objectif est
convexe mais \textbf{ni lisse, ni lineaire} ?
\end{center}

Ce cas de figure est loin d'etre marginal : il apparait des que l'on ajoute a
un objectif lisse un terme de regularisation non differentiable, situation
omnipresente en traitement du signal, en statistique et en apprentissage
automatique.

\subsection{La technique etudiee : les methodes de gradient proximal}

Pour repondre a cette question, nous etudions une famille de methodes de
premier ordre adaptees au cas non lisse : les \textbf{methodes de gradient
proximal}, et en particulier l'algorithme \textbf{ISTA} (\emph{Iterative
Shrinkage-Thresholding Algorithm}) et son acceleration \textbf{FISTA}
(\emph{Fast ISTA}). Conformement au sujet B, le c\oe ur de ce rapport est
l'\emph{etude de cette technique de resolution} : sa construction, ses garanties
de convergence, et sa comparaison avec les alternatives existantes.

Il importe de souligner d'emblee qu'ISTA et FISTA \emph{ne sont pas des
algorithmes d'apprentissage automatique} : ce sont des algorithmes
d'optimisation convexe, au meme titre que le gradient conjugue ou les methodes
de point interieur cites dans l'enonce du sujet B. Le fait qu'ils soient
frequemment employes en apprentissage ne change rien a leur nature : ils
resolvent un probleme d'optimisation, point.

\subsection{Fil conducteur et application illustrative}

Pour ancrer l'etude dans un probleme concret, nous utilisons comme fil rouge le
probleme du \textbf{LASSO} (regression lineaire avec penalisation $L_1$), qui
est l'instance canonique d'un objectif convexe composite \guillemotleft{}~lisse $+$ non
lisse~\guillemotright{}. En fin de rapport, nous illustrons l'interet pratique de la methode
sur une \textbf{application a la separation de sources audio}
(section~\ref{sec:audio}) ; cette application reste toutefois une
\emph{illustration} et non le sujet : elle rend le probleme tangible et fournit
une evaluation quantitative, mais l'analyse demeure centree sur l'optimisation.

\subsection{Organisation du rapport}

Le rapport s'organise en trois temps. La premiere partie pose le cadre :
rappels du cours mobilises (section~\ref{sec:rappels}) puis formulation precise
du probleme composite et preuve de sa convexite (section~\ref{sec:probleme}).
La deuxieme partie constitue le c\oe ur theorique : les outils de l'optimisation
non lisse, dont l'operateur proximal (section~\ref{sec:theorie}), puis les
algorithmes ISTA (section~\ref{sec:ista}) et FISTA (section~\ref{sec:fista})
avec leurs preuves de convergence. La troisieme partie met la methode en
perspective : etat de l'art comparatif (section~\ref{sec:etatart}), application
audio (section~\ref{sec:audio}) et experiences numeriques
(section~\ref{sec:exp}). La section~\ref{sec:conclusion} conclut.
